%% kapa: exam=1 round=22 year=2011 period=2 form=A subject=appraisal_law subject_ko=부동산관계법규 seq=30 number=70
%% source: Q-Net 감정평가사 기출문제 (cst00309) artlSeq=5036790 file=extracted/1008/5036790/2011년 제22회 감정평가사 1차시험 기출문제/2교시.hwp
%% conversion: kordoc 4.24.51 -> tools/gen1.py (auto); review=proofread-gpt-astra; flags=none
\begin{question}{70}
건축법령상 공개 공지의 확보 등에 관한 설명으로 옳지 않은 것은?
\begin{choices}
  \item 공개 공지의 면적은 대지면적의 100분의 10 이하의 범위에서 건축조례로 정한다.
  \item 문화 및 집회시설로서 연면적의 합계가 3천 제곱미터 이상일 경우 공개 공지를 확보하여야 한다.
  \item 공개 공지를 확보하는 경우 용적률은 해당 지역에 적용하는 용적률의 1.2배 이하의 범위 내에서 건축조례로 정하여 완화할 수 있다.
  \item 공개 공지에는 연간 60일 이내의 기간 동안 건축조례로 정하는 바에 따라 주민을 위한 문화행사를 열거나 판촉활동을 할 수 있다.
  \item 공개 공지에는 환경친화적으로 편리하게 이용할 수 있도록 건축조례로 정하는 시설을 설치하여야 한다.
\end{choices}
\end{question}
\answer{2}
